\section{Въведение в линейната регресия}

В тази лекция ще се запознаем с линейната регресия. Тя е един от най-основните инструменти в статистиката. Линейната регресия стои зад много съвременни инструменти и технологии, въпреки че често пъти е скрита вътре в самия модел, като например в невронните мрежи, рейтинговите системи и други.

Разбира се, както при всеки модел, ако предпоставките или допусканията не са изпълнени, лесно можете да стигнете до грешни изводи. Създаването на модели е изкуство, което изисква познаване на теорията и се овладява най-вече чрез практиката.

\subsection{Исторически контекст: Регрес към средното}
\begin{example}[Регрес към средното при ръста]\label{ex:14-1}
Началото на линейната регресия е свързано с Франсис Галтън, който се опитал да изучи зависимостта между ръста на бащите и ръста на техните синове. В своето изследване Галтън се опитал да предскаже височината на сина като функция на ръста на бащата. Забелязал е, че един добър модел може да се запише във вида:
\[
\text{Ръст на сина} = \text{Среден ръст} + \beta (\text{Ръст на бащата} - \text{Среден ръст}) + \varepsilon
\]
В изчисленията на Галтън коефициентът $\beta$ е бил приблизително $0{,}6$. Ако бащата е с 10 см по-висок от средния ръст за популацията, моделът предвижда отклонение
\[
0{,}6\cdot 10\text{ см}=6\text{ см}
\]
за сина: средно около 6 см над средния ръст, а не цели 10 см.

Това явление се нарича \textbf{регрес към средното}. То се забелязва нерядко и при други човешки характеристики, като например интелекта. Изключително умни хора имат тенденцията да имат наследници, чиито показатели са по-близо до средното ниво на популацията. Името „регресия“ произхожда именно от това „връщане“ към средната стойност.
\end{example}

Най-общо, един такъв модел може да се запише като:
\[
Y_k = \beta_0 + \beta_1 X_k + \varepsilon_k
\]
където $X_k$ е предикторът (например ръстът на бащата или дозата на дадено лекарство), $Y_k$ е откликът (ръстът на сина или реакцията към лекарството), $\beta_0$ и $\beta_1$ са линейни коефициенти, а $\varepsilon_k$ е стохастична грешка.

\section{Детерминистичен модел. Метод на най-малките квадрати}

Нека разгледаме детерминистичната оптимизационна задача, която стои зад линейната регресия. Имаме дадени $n$ на брой наблюдения, представляващи точки в равнината $(X_k, Y_k)$ за $k=1, \dots, n$. Търсим права линия с уравнение $y_k = b_0 + b_1 X_k$, която да моделира данните по възможно най-добрия начин. Грешката за дадено наблюдение $k$ е разликата между истинската стойност $Y_k$ и предвидената стойност от нашата права $\hat{Y}_k = b_0 + b_1 X_k$.

\begin{defn}[метод на най-малките квадрати]\label{def:ols}
\emph{Оценки по метода на най-малките квадрати} наричаме онези стойности на $b_0$ и $b_1$, които минимизират сумата от квадратите на грешките:
\[
\sum_{k=1}^n \varepsilon_k^2 = \sum_{k=1}^n (Y_k - b_0 - b_1 X_k)^2 \to \min .
\]
\end{defn}

Изборът на квадрата на грешката (вместо модул) исторически е бил мотивиран от изчислително удобство, тъй като функцията е гладка и лесно се диференцира, а също така съответства на Евклидовото разстояние. Освен това, от статистическа гледна точка, този подход води до добри оценки.

\subsection{Намиране на коефициентите}
За да минимизираме функцията, намираме нейните частни производни спрямо $b_0$ и $b_1$ и ги приравняваме на нула. Ще използваме стандартните означения $\bar{X} = \frac{1}{n}\sum X_k$ и $\bar{Y} = \frac{1}{n}\sum Y_k$ за средните аритметични стойности.

Диференцирайки по $b_0$ и разделяйки на $-2$, получаваме първото уравнение (нормално уравнение):
\[
\sum_{k=1}^n (Y_k - b_0 - b_1 X_k) = 0
\]
\[
\sum_{k=1}^n Y_k - n b_0 - b_1 \sum_{k=1}^n X_k = 0
\]
Изразявайки сумите чрез средните аритметични ($n \bar{Y} - n b_0 - b_1 n \bar{X} = 0$), след съкращаване на $n$ стигаме до връзката за $b_0$:
\[
b_0 = \bar{Y} - b_1 \bar{X}
\]

Диференцирайки по $b_1$, получаваме второто уравнение:
\[
\sum_{k=1}^n X_k (Y_k - b_0 - b_1 X_k) = 0
\]
\[
\sum_{k=1}^n X_k Y_k - b_0 n \bar{X} - b_1 \sum_{k=1}^n X_k^2 = 0
\]

Ако заместим $b_0$ във второто уравнение и извършим алгебричните преобразувания, можем да изведем явна формула за наклона $b_1$. Формулата може да се запише в няколко еквивалентни и много полезни вида:
\[
b_1 = \frac{\sum_{k=1}^n X_k Y_k - n \bar{X}\bar{Y}}{\sum_{k=1}^n X_k^2 - n \bar{X}^2} = \frac{\sum_{k=1}^n (X_k - \bar{X})(Y_k - \bar{Y})}{\sum_{k=1}^n (X_k - \bar{X})^2}
\]
Забележете, че горните изрази изключително много приличат на формулата за ковариация и дисперсия, откъдето идва и връзката с коефициента на корелация.

\subsection{Удобен запис за статистическия анализ}
За да изследваме свойствата на $b_0$ и $b_1$ по-лесно, ще въведем следното означение. Нека $A$ е сумата от квадратите на отклоненията на предикторите от тяхното средно:
\[
A = \sum_{k=1}^n (X_k - \bar{X})^2
\]
Въвеждаме и теглата $c_k$:
\[
c_k = \frac{X_k - \bar{X}}{A}
\]
Тогава оценката $b_1$ може елегантно да се запише като линейна комбинация на отклиците $Y_k$:
\[
\hat{\beta}_1 = b_1 = \sum_{k=1}^n c_k Y_k
\]
Оценката $\hat{\beta}_0$ може да се запише по аналогичен начин, замествайки $b_1$ в $b_0 = \bar{Y} - b_1 \bar{X}$:
\[
\hat{\beta}_0 = b_0 = \sum_{k=1}^n \left( \frac{1}{n} - \bar{X} c_k \right) Y_k
\]
Много важно свойство на константите $c_k$, което ще използваме често, е, че сумата им е нула:
\[
\sum_{k=1}^n c_k = \frac{\sum_{k=1}^n (X_k - \bar{X})}{A} = \frac{n\bar{X} - n\bar{X}}{A} = 0
\]

\section{Стохастичен модел}

Дотук разгледахме задачата детерминистично. По-общо, моделът се разглежда статистически. Предикторите $X_k$ се считат за детерминистични величини (например дозата от дадено лекарство, която ние контролираме и избираме), докато отклиците $Y_k$ са случайни величини. Стохастиката в $Y_k$ идва единствено от случайната грешка $\varepsilon_k$.

\begin{defn}[стохастичен модел на простата линейна регресия]\label{def:slr}
Истинският модел е
\[
Y_k = \beta_0 + \beta_1 X_k + \varepsilon_k ,
\]
при следните допускания за грешките $\varepsilon_k$:
\begin{enumerate}
    \item \textbf{Независимост:} случайните величини $(\varepsilon_1, \dots, \varepsilon_n)$ са взаимно независими.
    \item \textbf{Липса на систематична грешка:} $\E[\varepsilon_k] = 0$ за всяко $k$, откъдето $\E[Y_k] = \beta_0 + \beta_1 X_k$.
    \item \textbf{Хомоскедастичност (еднаква дисперсия):} $\Var(\varepsilon_k) = \sigma^2$ за всяко $k$.
    \item \textbf{Нормалност:} грешките са нормално разпределени, $\varepsilon_k \sim N(0, \sigma^2)$.
\end{enumerate}
\end{defn}

Тези допускания са силни. Второто означава, че средно моделът е точен; третото — че грешката има еднаква вариативност, независимо от стойността на предиктора $X_k$.

\textbf{Забележка относно допусканията:} Обобщения на този модел има в най-различни посоки (напр. хетероскедастичност). Винаги обаче трябва да наложите някаква структура на грешката, за да можете да правите изводи. Ако човек просто използва модела без да проверява дали тези допускания са валидни (както се е случвало често в практиката преди финансовата криза от 2008-2009 г.), той може да получи напълно погрешни и катастрофални резултати, защото моделът ще върне число, но то може да е напълно невалидно.

При допусканията на модела от дефиниция~\ref{def:slr} се оказва, че оценките $\hat{\beta}_0 = b_0$ и $\hat{\beta}_1 = b_1$, получени по метода от дефиниция~\ref{def:ols}, съвпадат с \textbf{максимално правдоподобните оценки} (MLE).

\section{Свойства на оценките}

Използвайки допусканията, ще изследваме свойствата на нашите оценки. Тъй като $b_1$ и $b_0$ са линейни комбинации на нормално разпределени случайни величини ($Y_k$), то самите те също ще имат нормално разпределение. Трябва само да намерим тяхното математическо очакване и дисперсия.

\subsection{Оценка на наклона $\hat{\beta}_1$ ($b_1$)}
Проверяваме дали оценката $b_1$ е неизместена (unbiased). Математическото очакване е:
\[
\E[b_1] = \E\left[ \sum_{k=1}^n c_k Y_k \right] = \sum_{k=1}^n c_k \E[Y_k] = \sum_{k=1}^n c_k (\beta_0 + \beta_1 X_k)
\]
Тъй като $\sum_{k=1}^n c_k = 0$, членът с $\beta_0$ отпада:
\[
\E[b_1] = \beta_1 \sum_{k=1}^n c_k X_k
\]
Сега нека пресметнем сумата $\sum c_k X_k$. Тъй като можем да добавим $-\bar{X}$ (което е константа) без да променим стойността на сумата (защото $\bar{X} \sum c_k = 0$):
\[
\sum_{k=1}^n c_k X_k = \sum_{k=1}^n \frac{X_k - \bar{X}}{A} X_k = \sum_{k=1}^n \frac{X_k - \bar{X}}{A} (X_k - \bar{X}) = \frac{\sum (X_k - \bar{X})^2}{A} = \frac{A}{A} = 1
\]
Следователно $\E[b_1] = \beta_1$, което доказва, че оценката е \textbf{неизместена}.

Дисперсията на $b_1$ намираме използвайки факта, че $Y_k$ са независими (тъй като $\varepsilon_k$ са независими) и всяко има дисперсия $\sigma^2$:
\[
\Var(b_1) = \Var\left( \sum_{k=1}^n c_k Y_k \right) = \sum_{k=1}^n \Var(c_k Y_k) = \sum_{k=1}^n c_k^2 \Var(Y_k) = \sum_{k=1}^n c_k^2 \sigma^2
\]
Заместваме $c_k$:
\[
\Var(b_1) = \sigma^2 \sum_{k=1}^n \left( \frac{X_k - \bar{X}}{A} \right)^2 = \sigma^2 \frac{\sum (X_k - \bar{X})^2}{A^2} = \frac{\sigma^2 A}{A^2} = \frac{\sigma^2}{A}
\]
В крайна сметка получаваме, че $b_1$ има следното нормално разпределение:
\[
b_1 \sim N\left(\beta_1, \frac{\sigma^2}{A}\right)
\]

\subsection{Оценка на пресечната точка $\hat{\beta}_0$ ($b_0$)}
Аналогично пресмятаме очакването на $b_0$:
\[
\E[b_0] = \sum_{k=1}^n \left( \frac{1}{n} - \bar{X} c_k \right) \E[Y_k] = \sum_{k=1}^n \left( \frac{1}{n} - \bar{X} c_k \right) (\beta_0 + \beta_1 X_k)
\]
Ако се разкрият скобите и се използват вече доказаните свойства ($\sum c_k = 0$ и $\sum c_k X_k = 1$), се получава, че $\E[b_0] = \beta_0$. Оценката $b_0$ също е неизместена.

Дисперсията на $b_0$ е:
\[
\Var(b_0) = \sum_{k=1}^n \left( \frac{1}{n} - \bar{X} c_k \right)^2 \Var(Y_k) = \sigma^2 \sum_{k=1}^n \left( \frac{1}{n} - \bar{X} c_k \right)^2
\]
След алгебрични преработки, които ще спестим за краткост, се получава:
\[
\Var(b_0) = \sigma^2 \left( \frac{1}{n} + \frac{\bar{X}^2}{A} \right)
\]
Следователно, разпределението на $b_0$ е:
\[
b_0 \sim N\left(\beta_0, \sigma^2 \left(\frac{1}{n} + \frac{\bar{X}^2}{A}\right)\right)
\]

\section{Оценяване на дисперсията $\sigma^2$}

Ако дисперсията $\sigma^2$ на грешките ни е предварително известна, то тогава можем веднага да използваме горните разпределения. Често пъти обаче $\sigma^2$ не е известна и следва да я оценим от данните.

Нека разгледаме стандартизираните грешки:
\[
Z_k = \frac{Y_k - \beta_0 - \beta_1 X_k}{\sigma} \sim N(0, 1)
\]
Тъй като тези величини са независими стандартни нормални, сумата от техните квадрати има хи-квадрат разпределение с $n$ степени на свобода:
\[
\sum_{k=1}^n Z_k^2 = \frac{1}{\sigma^2} \sum_{k=1}^n (Y_k - \beta_0 - \beta_1 X_k)^2 \sim \chi^2(n)
\]
Проблемът е, че ние не знаем истинските параметри $\beta_0$ и $\beta_1$. Можем да ги заместим с техните оценки $b_0$ и $b_1$, които сме пресметнали. Тъй като използваме две оценки (губим две степени на свобода), разпределението се променя на $\chi^2(n-2)$:
\[
\frac{1}{\sigma^2} \sum_{k=1}^n (Y_k - b_0 - b_1 X_k)^2 \sim \chi^2(n-2)
\]

Според Закона за големите числа (ЗГЧ), ако разделим случайна величина с $\chi^2(n-2)$ разпределение на нейните степени на свобода, то при $n \to \infty$ тази дроб клони към 1 (математическото очакване на $Z_1^2$). Това ни подсказва, че можем да конструираме \textbf{силна асимптотично състоятелна оценка} за дисперсията $\sigma^2$:
\[
\hat{\sigma}^2 = \frac{1}{n-2} \sum_{k=1}^n (Y_k - b_0 - b_1 X_k)^2
\]
В ситуации, в които $\sigma^2$ е неизвестна, просто я заместваме с $\hat{\sigma}^2$ във всички по-нататъшни разсъждения.

\section{Проверка на хипотези}

Често се налага да проверяваме хипотези относно линейните коефициенти. Най-често нулевата хипотеза е, че предикторът не влияе на отклика, т.е. че наклонът $\beta_1$ е нула: $H_0: \beta_1 = 0$. По-общо, можем да тестваме дали параметърът приема конкретна стойност $H_0: \beta_1 = \beta$, срещу алтернативата $H_1: \beta_1 \neq \beta$.

\subsection{Случай на известно $\sigma^2$}
Когато $\sigma^2$ е известно, използваме факта, че $b_1 \sim N(\beta_1, \sigma^2/A)$. При валидност на нулевата хипотеза $H_0: \beta_1 = \beta$, статистиката $Z$ е стандартно нормално разпределена:
\[
Z = \frac{b_1 - \beta}{\sqrt{\frac{\sigma^2}{A}}} \sim N(0, 1)
\]
Ако изберем ниво на значимост (грешка от първи род) $\alpha$, критичната област $W$, в която отхвърляме нулевата хипотеза, е симетрична:
\[
W = \left\{ |Z| \geq q_{1-\frac{\alpha}{2}} \right\}
\]
където $q_{1-\frac{\alpha}{2}}$ е квантилът на стандартното нормално разпределение.
Оттук лесно можем да построим и доверителен интервал за $\beta_1$ с вероятност $1-\alpha$:
\[
\beta_1 \in \left( b_1 - q_{1-\frac{\alpha}{2}}\sqrt{\frac{\sigma^2}{A}} \; , \; b_1 + q_{1-\frac{\alpha}{2}}\sqrt{\frac{\sigma^2}{A}} \right)
\]

\subsection{Случай на неизвестно $\sigma^2$}
Когато $\sigma^2$ не е известно, ние заместваме $\sigma^2$ с нейната оценка $\hat{\sigma}^2$. Разпределението на тестовата статистика се променя от нормално на \textbf{t-разпределение на Стюдънт} с $n-2$ степени на свобода:
\[
T = \frac{b_1 - \beta}{\sqrt{\frac{\hat{\sigma}^2}{A}}} \sim t(n-2)
\]
При конструиране на критична област и доверителни интервали, вместо квантилите на нормалното разпределение $q_{1-\alpha/2}$, трябва да използваме квантилите на $t$-разпределението $t_{1-\alpha/2, n-2}$, които имат по-тежки опашки.

\textbf{Важна бележка:} Когато $n$ е достатъчно голямо (обикновено $n \geq 32$, така че $n-2 \geq 30$), $t$-разпределението е много близко до стандартното нормално разпределение $N(0,1)$. В такива случаи можем безопасно да използваме нормалните квантили $q_{1-\alpha/2}$ като апроксимация, третирайки статистиката като нормална, просто замествайки $\sigma^2$ с $\hat{\sigma}^2$.

\subsection{Тестване на пресечната точка $\beta_0$}
Абсолютно същата логика може да се приложи и за тестване на хипотеза за пресечната точка: $H_0: \beta_0 = \beta^*$.
При известно $\sigma^2$ статистиката е:
\[
Z = \frac{b_0 - \beta^*}{\sqrt{\sigma^2\left(\frac{1}{n} + \frac{\bar{X}^2}{A}\right)}} \sim N(0, 1)
\]
Ако $\sigma^2$ не е известно, отново използваме $T$-статистика с $\hat{\sigma}^2$ и $t$-разпределение.
